\section{Implementation walkthrough and technical discussion}
\label{implementation-walkthrough}
This chapter connects the product behavior to the code that implements it.
The following chapters provide the full pack schema, verifier contracts, API,
commands, storage migrations and deployment boundaries. Read the walkthrough
first, then use those references to inspect a specific implementation. Descriptions
refer to implemented behavior; limitations and possible extensions are identified
explicitly rather than presented as shipped features.

\subsection{A useful way to explain Ascend}
Ascend is a self-hosted practice tracker. Authored TOML packs describe a dependency
graph of tasks and levels. A learner selects a mission, performs an eligible task,
records work and, where supported, submits evidence. The application keeps learning
notes, progression, streaks and verified rewards distinct. Reading a resource,
installing a pack or letting a timer run does not prove completion.

The Python engine expresses rules, SQLite stores account-scoped state, and
\texttt{Tracker} coordinates use cases. The CLI and FastAPI share those operations.
The browser is a small vanilla JavaScript client with locally served assets, not
a separate implementation of scoring or progression. Native exercises compare
constrained submitted data with server-held expected results. External verifier
adapters have separate trust and operational boundaries.

\subsection{Code map and responsibilities}
\begin{longtable}{P{.25\textwidth}P{.67\textwidth}}
\textbf{Area} & \textbf{Responsibility and where to read}\\\hline
Content and registry & \path{packs/community/} contains canonical TOML by domain.
\path{ascend/registry.py} resolves stable IDs to registry paths;
\path{scripts/build_registry.py} rebuilds registry metadata.\\
Pure rules & \path{ascend/engine/packs.py}, \path{ascend/engine/levels.py},
\path{ascend/engine/streak.py} and \path{ascend/engine/validation.py} parse content,
validate constraints and derive eligibility or streak results.\\
Application use cases & \path{ascend/application.py} owns installation, missions,
entries, marks, evidence, undo and derived overview/statistics. Follow
\texttt{Tracker.log}, \texttt{check}, \texttt{levels} and \texttt{overview}.\\
Storage & \path{ascend/storage/__init__.py} supplies scoped access and transactions.
\path{ascend/storage/migrations/} is the ordered schema history.\\
Sessions and grading & \path{ascend/timing.py} owns server time and XP;
\path{ascend/labs/__init__.py} generates fixtures, returns public views and grades
typed answers. \path{ascend/labs/cases.json} supplies reasoning cases.\\
Transports & \path{ascend/cli/__init__.py} parses and dispatches commands.
\path{ascend/cli/remote.py} connects the CLI to a server.
\path{ascend/server/app.py} authenticates and exposes HTTP operations.\\
Trust and identity & \path{ascend/server/auth.py} owns passwords and bearer tokens.
\path{ascend/verifiers/} owns file, Git, HTTP, command, plugin and CTF adapters.\\
Reminders & \path{ascend/nudge/} decides when to notify and supplies channel adapters.
\path{ascend/server/telegram.py} processes authenticated reminder interactions.\\
Browser state & \path{web/app.js} owns routing, rendering, forms, feedback and
server-synchronized session display. \path{web/index.html} provides semantic regions.\\
Presentation & \path{web/style.css} owns layout and components.
\path{web/effects.css} and \path{web/effects.js} add optional motion;
\path{web/theme.js} applies saved appearance before first paint.\\
Offline help & \path{web/sw.js} caches an explicit application-shell allowlist.
\path{web/guide.json} is the source for both in-app help and a generated PDF chapter.\\
Verification & \path{tests/} exercises rules, transports and browser behavior.
\path{scripts/ui_screenshots.py} captures disposable accounts and optional axe checks;
\path{scripts/build_docs.py} generates references and checks PDF compilation.\\
\end{longtable}

\subsection{Follow a pack from repository to recommendation}
An installed pack has a stable ID independent of its folder or display title.
Registry entries carry relative paths and checksums. The domain folders make the
repository navigable without changing the IDs used by learner state. This is why
moving a file must update registry generation and packaging rather than rename
task identities. The ethical-hacking foundations curriculum is separate from
offensive specialization; the introductory track does not stand in for every
offensive role or skill.

Installation parses and validates TOML before storing a private copy and metadata.
The pack graph defines prerequisites, levels and milestone gates. The application
combines those authored rules with account-specific marks and explicit overrides
to derive candidates. Today selects an eligible next task; its time filter changes
the recommendation, not the underlying task state. A quick win is an alternative,
not a second mandatory assignment.

The browser's mission picker first chooses a domain, then an installed track in
that domain. Known installed IDs are reconciled against the catalog when an older
server response omits category metadata or reports Custom. Unknown private packs
remain custom. Changing domains clears the track selection. Installing content,
starting a mission, starting a timed task and completing a task are four separate
operations, each with its own visible action.

\subsection{Six records with different responsibilities}
\begin{longtable}{P{.23\textwidth}P{.69\textwidth}}
\textbf{Record} & \textbf{Meaning and invariant}\\\hline
Mission & The selected learning track: active, paused or done. At most one active
mission per account. Mission pausing does not stop a task timer.\\
Node state & A task is todo, done or skipped. Eligibility is derived from this state
and the graph; a running timer alone does not change it.\\
Task session & One scored timing history per account/pack/node. Current actions
create a running session or finish it. A unique partial index limits running
sessions to one per account. Legacy paused rows are retained, not silently rewritten.\\
Exercise instance & Randomized fixture, private expected answer, attempt count,
fingerprint and pass timestamp. A fresh instance replaces the current fixture;
it does not create a new scoring history.\\
XP award & One immutable reward per account/pack/node, backed by a database primary
key and an idempotent insert. An award is not a mutable browser counter.\\
Learning entry & A dated account-scoped note with optional task/proof/minutes and
completion metadata. Entries drive streak and time summaries; a task mark alone
does not necessarily create a learning day.\\
\end{longtable}

\subsection{Trace Start, reload, verify and End}
The challenge preview is a dialog containing the finish condition, steps, resource,
hints and session target. Pressing Start moves into a full page with a stable hash
route, such as \texttt{\#challenge/PACK:NODE}. The colon is URL-encoded when the
application creates that route. The preview dialog closes; exercise data and answer
fields use the shared page width rather than a nested modal.

For a bundled exercise, the browser sends \texttt{lab start --minutes N} through
\texttt{POST /api/command}. \texttt{labs.start} opens a transaction and calls
\texttt{timing.start} inside it. Nested transactions reuse the same connection's
transaction. Prerequisites, target bounds and the one-running-session rule are
checked before the fixture is committed. A failed setup cannot leave a newly
created timer without its exercise. An ordinary task uses \texttt{session start}.

The browser then reads \path{GET /api/challenges/{reference}}. That endpoint resolves
the installed pack independently of the currently selected mission and returns:
\begin{verbatim}
{
  "node": {"id": "...", "pack_id": "...", "status": "...", "...": "..."},
  "level": {"id": "...", "title": "...", "open": true},
  "session": {"status": "running", "...": "..."},
  "exercise": {"id": "...", "data": {}, "answer_template": {}, "...": "..."},
  "exercise_error": null
}
\end{verbatim}
This is a shape illustration; optional session/exercise values can be null and
the authoritative fields come from the code. Expected answers are never included.
Reading uses \texttt{labs.current} and \texttt{view}, not \texttt{start}. It therefore
does not create a fixture or restart a clock. Invalid, expired or changed instances
produce an exercise error while preserving access to the authored task.

A reload reconstructs the workspace from the route and this read endpoint. The
session strip on every page offers Resume challenge, which navigates to the same
reference. Draft answers use session storage keyed by account name and instance ID;
they survive reloads in that tab, are isolated from another instance, and are removed
on successful verification or an explicit fresh instance. They are not a cross-device
draft service. Emitting a logging event updates the displayed data without rebuilding
the answer form or losing typed values.

Submission sends typed answers and an instance ID, not a score. The server checks
ownership, latest-instance identity, pack fingerprint, expiry and attempt limits.
Exact keys and types must match: Python's ordinary equality would otherwise accept
\texttt{True == 1}. A passing result records verification and invokes the shared
completion path. Repeated submission of a passed instance is idempotent. Invalid
answers consume an attempt but grant neither completion nor XP.

End session is an explicit alternative to completion. It stores elapsed time and a
finished timestamp, without marking the node done or granting XP. A later successful
verification may demonstrate the task but cannot resurrect the ended attempt's
reward. The clock cannot be paused. Old paused rows are readable for unscored
practice and may be ended, but cannot be resumed or used to gain a reward. The
separate control named Pause motion effects changes decoration only.

\subsection{Why the timer is server authoritative}
\texttt{running\_since}, the initial budget and accumulated seconds live in SQLite.
While running, elapsed time is calculated from the server clock. The browser stores
the last server summary and its local \texttt{performance.now()} reading, then
interpolates a display each second. Visible clients refresh the summary periodically.
The local countdown is a convenience; changing JavaScript, suspending the tab or
editing displayed text cannot change server grading time.

Reopening the same task preserves its original budget and start time. Database
constraints prevent concurrent duplicate sessions and awards. The transaction
boundary is necessary in addition to disabled buttons: two tabs or clients can
submit requests concurrently. \texttt{BEGIN IMMEDIATE}, uniqueness constraints and
idempotent completion protect persisted invariants; a spinner protects only the
interaction in one view. There is no claim that learner-selected targets or
administrator-controlled clocks constitute a proctored exam.

\subsection{Browser navigation, layout and feedback}
\texttt{currentRoute} parses the hash; \texttt{switchTab} updates the visible section,
navigation state and document title before rendering the destination. Main screens
are Today, Path, Stats, Tracks, Docs and Profile; Challenge is a task workspace reached
from a task or session. A refresh-generation counter prevents an older overview
response from replacing newer data. Challenge rendering also checks that the user
is still on the requested reference after awaiting its response. This is targeted
stale-response protection, not a general request-cancellation framework.

All page frames and the top bar share a maximum width. Dashboard support follows
the main challenge as open rows, Stats uses aligned metrics and plots, and settings
use sections and dividers. Cards remain where a discrete summary needs emphasis.
Grid children use shrinkable columns and minimum-width rules; bounded code/data
regions scroll internally and are keyboard focusable. Breakpoint tests check
containment and horizontal overflow rather than merely checking CSS declarations.

Docs separates navigation from reading. The desktop rail groups ten topics; compact
screens move the same navigation DOM into a native dialog drawer. Search filters
inside navigation, leaving article geometry unchanged. Native dialog modality,
explicit initial focus, Escape handling and focus restoration support keyboard use.
The article title and body share an alignment edge and topic links survive reloads.

\texttt{runAction} prevents duplicate activation while a request is pending and
restores the button on completion. User-controlled and authored strings become
DOM text, not injected HTML. Feedback is moved into the active dialog's feedback
region so it is not hidden behind the browser's top layer. Errors remain visible
until dismissed or replaced. Timer reminders do not replace unresolved errors.

\subsection{Visual effects without a framework migration}
The visual research considered React Bits, Aceternity UI, Magic UI and Motion's
accessibility guidance. Their patterns were adapted in original CSS and browser
animation code; the app does not import React, GSAP, Motion or a WebGL renderer.
The source register records which rendered pages or repository sources were read.

The ambient field combines violet/teal radial gradients, a masked dot pattern,
sparse particles and an SVG orbit. Slow transform animations give depth while
Docs lowers their visual intensity. \texttt{AscendMotion.enter} performs brief
heading and section entrances and animates Stats bars; text stays complete in the
DOM. A route key prevents data refreshes from replaying an entrance. The separately
guarded next-challenge border sweep plays once per page load.

Hover uses normal pointer and keyboard affordances. A small radial highlight follows
the native pointer only over the main challenge on devices with a fine pointer.
It ignores pointer events, stays inside the card and is disabled with motion. There
is no custom cursor, scroll interception or continuous JavaScript render loop.
Background animation uses CSS; heading transitions use the Web Animations API.
These choices bound the work, but do not establish a measured frame-rate guarantee.

A persistent top-bar control pauses decorative effects. Profile exposes the saved
Calm/Off preference. The system's reduced-motion preference always wins, including
changes while the page is open; pending entrance animations are cancelled.
Background animations pause when the document is hidden. Reduced motion retains
readable static content, semantic status, focus rings and the server timer.

\subsection{Persistence, offline behavior and trust boundaries}
SQLite runs with WAL, foreign keys and a bounded busy timeout. Request dependencies
open an account-scoped database handle and close it afterwards. Migration files
are applied in order inside transactions; the schema version records progress.
The generated storage chapter reproduces the real SQL, including uniqueness and
partial-index constraints. File copies, exercise state, node marks, notes and XP
must be backed up together for a recoverable profile.

The service worker caches only an explicit shell list: local scripts, styles, guide,
font and icons. It tries the network first and falls back to that cache. API responses
are no-store and excluded from this handler. Offline help is supported; offline
progress writes, grading and a cross-device sync queue are not. Adding a shell asset
requires updating the cache list and cache generation, plus verifying the built wheel.

Password hashing, token storage, account scoping, content fingerprints and executable
trust solve different problems. Scrypt protects password verification; hashed bearer
tokens authenticate requests; scoped queries isolate records; fingerprints revoke
stale trust after content changes. Plugins run with the host user's authority despite
process isolation. HMAC external observations trust the administrator's collector and
key. None of these mechanisms establishes who physically solved a shared external lab.

\subsection{Validation and debugging workflow}
Use a disposable profile to reproduce a report. Identify whether the problem is
authored content, derived eligibility, persisted state, transport or presentation.
Inspect the API response before changing a rendered label. A timer-without-exercise
report requires checking the route and saved instance identity, not merely adding
another timer display. A layout report requires measuring visible bounding boxes
at the reported width, then inspecting screenshots and interaction states.

The automated suite combines pure rule tests, storage/transaction tests, CLI and HTTP
tests, real-browser flows and a wheel installed outside the source checkout. Browser
regressions cover reload recovery, answer drafts, preserved timing, end-without-XP,
navigation, keyboard behavior, offline guide assets and reduced motion. The screenshot
script adds geometry and selected axe rules. Passing these checks does not replace
human usability review, assistive-technology testing or deployment-specific testing.

For a feature change, update the shared use case first, adapt its transports and UI,
add a regression that would fail for the original bug, then refresh generated docs.
Check formatting, selected lint rules, JavaScript syntax, browser behavior and the
affected PDFs. Deploy using the same profile and environment after a SQLite backup;
compare protected records across a restart. Never create demonstration progress in
the owner's account. Keep historical evidence labeled with the build it tested.

\subsection{Interview preparation: explain the implementation honestly}
\begin{enumerate}
\item \textbf{Why separate the engine, application and transports?}
Pure rules are testable without HTTP or a browser; \texttt{Tracker} owns shared use
cases so local CLI and server requests do not drift into different business rules.
\item \textbf{How is progression calculated?}
The pack dependency graph, level gates, account node states and explicit overrides
produce eligible candidates. The UI displays these results rather than awarding
unlocks based on its own click history.
\item \textbf{How do you prevent duplicate XP?}
The account/pack/node primary key, transactional completion and idempotent inserts
protect the invariant across retries and concurrent clients. Disabling a button
is useful feedback but is insufficient as a data-integrity mechanism.
\item \textbf{Why did reload recovery need a read endpoint?}
The route needs enough identity to fetch the existing instance. Reusing a start
operation for a read could create new random data or alter timing; recovery must
be side-effect free with respect to the exercise and session.
\item \textbf{What happens if exercise setup fails?}
Timer creation and fixture creation share a transaction. An exception rolls them
back together. Errors remain visible and the user can retry without an orphaned
new timer being mistaken for a successfully opened exercise.
\item \textbf{Can a user stop the clock and submit later for more XP?}
No. End finishes timing without reward; only a previously running session can
receive qualifying XP during completion. A fresh fixture cannot reset that attempt.
\item \textbf{Why distinguish a mission pause from a task session?}
A mission expresses the selected track. A task session is a timed attempt. Changing
the mission does not stop elapsed time, and task pausing is no longer an action.
\item \textbf{What is actually verified?}
A native fixture verifies a constrained result; a reasoning case verifies a choice;
a signed observation verifies a collector's report; a file check verifies limited
artifact properties. None alone proves broad professional competence.
\item \textbf{How are accounts isolated?}
Authentication selects an account-scoped tracker, and storage helpers scope records.
Exercise reads and submissions check ownership and current-instance identity before
returning public data or accepting answers.
\item \textbf{What does the PWA support offline?}
The installed shell and help can load. Private API data is not cached, and there is
no offline write queue. State-changing actions still require the server.
\item \textbf{Why use an overlay menu for compact Docs?}
It preserves reading position and article width. The same topic list moves into a
native modal with controlled focus instead of expanding above the article.
\item \textbf{Why avoid importing an animation framework?}
The selected effects need only gradients, SVG, transforms and short element
animations. Existing browser APIs meet that scope without changing the app's stack.
The choice can be reconsidered if future interaction complexity warrants it.
\item \textbf{How is accessibility checked?}
Semantic regions, labels, keyboard/focus behavior and reduced motion are implemented
directly, then tested with browser regressions and selected axe rules. Do not describe
an automated clean result as complete accessibility certification.
\item \textbf{What limits scalability?}
SQLite serializes writers; exercise and UI operations share one server process's
services and scheduler. WAL helps readers coexist with writes, but a multi-instance
scheduler or distributed write workload would need additional coordination. Such a
distributed redesign is not implemented here.
\item \textbf{What would you demonstrate live?}
Install a track in a disposable profile, start a task, reload into the same instance,
submit a wrong and then correct result, inspect one reward, and show that an ended
attempt cannot earn XP later. Also demonstrate keyboard Docs navigation and reduced
motion. Tie each observation to a test and source function.
\end{enumerate}
